\chapter{Einführung}
\label{kap:einfuehrung}

\section{Was ist ESIM?}

ESIM ist ein Simulator für analoge elektronische Schaltungen und für
differential-algebraische Gleichungssysteme (DAE). Er besteht aus
mehreren Teilen, die zusammenarbeiten, aber auch einzeln nutzbar sind:

\begin{description}[style=nextline,leftmargin=2em]
  \item[\esim\ – der Simulator]
    Ein Kommandozeilenprogramm in C++. Es liest eine Eingabedatei im
    JSON-Format, rechnet die darin verlangten Analysen und schreibt die
    Ergebnisse als CSV-Dateien. Der Arbeitspunkt und Messgrößen erscheinen
    zusätzlich auf der Konsole.
  \item[\editor\ – die Oberfläche]
    Ein Programm mit Fenstern (Qt 6) für macOS und Windows: Schaltplan und
    Text nebeneinander, Prüfung der Eingabe beim Tippen, Simulation per
    Tastendruck und ein Kurvenbetrachter mit Zoom, Fadenkreuz,
    Histogrammen und Kennlinienscharen.
  \item[Werkzeuge]
    \datei{tools/spice2json.py} übersetzt SPICE-Netzlisten,
    \datei{tools/plot.py} zeichnet Ergebnisse mit matplotlib, und
    \datei{python/esim.py} ruft den Simulator aus Python auf und liefert die
    Ergebnisse als Datenstrukturen.
\end{description}

\section{Was lässt sich simulieren?}

\begin{itemize}
  \item \textbf{Schaltungen} aus Widerständen, Kondensatoren, Spulen und
    gekoppelten Spulen, unabhängigen und gesteuerten Quellen,
    Verhaltensquellen mit beliebigen Ausdrücken, Dioden,
    Bipolartransistoren (Gummel-Poon), Sperrschicht-FETs, MOSFETs (Level 1,
    EKV, Level 3 und BSIM3v3), idealen Schaltern, idealen Dioden sowie
    verlustlosen und verlustbehafteten Leitungen.
  \item \textbf{Gleichungssysteme}, die direkt in der Eingabedatei stehen:
    gewöhnliche Differentialgleichungen mit algebraischen Nebenbedingungen,
    auch vom Index 3 und höher (sie werden automatisch umgeformt), Vektoren
    von Unbekannten, Totzeiten und Tabellen.
  \item \textbf{Ereignisse}: Umschaltungen und Sprünge, die vom Zustand
    oder von der Zeit abhängen – Schalter, Anschläge, Hysterese, abgetastete
    Regler.
  \item \textbf{Alles gekoppelt} in einem System, etwa eine Schaltung mit
    der Wärmebilanz eines Bauteils und einem digitalen Regler.
\end{itemize}

Die Analysen sind Arbeitspunkt, DC-Sweep, Frequenzgang (AC), Zeitverlauf
(Transient), Rauschen, periodischer stationärer Zustand, Fourier-Analyse,
Empfindlichkeiten und die Anpassung von Parametern an Messdaten. Jede
Analyse lässt sich über Parameter-Sweeps und Monte-Carlo-Läufe wiederholen
und mit Messgrößen (Grenzfrequenz, Anstiegszeit, Effektivwert …)
auswerten.

\section{Das Prinzip}

Der numerische Kern löst allgemein das System
\[
  \frac{d}{dt}\,q(x) + f(x,t) = 0 .
\]
Dabei ist $x$ der Vektor aller Unbekannten (Knotenspannungen, Zweigströme,
Variablen von Gleichungen). Zeilen, in denen $q$ verschwindet, sind
algebraische Bedingungen, alle anderen Differentialgleichungen. Eine
Schaltung liefert $f$ (Ströme) und $q$ (Ladungen und magnetische Flüsse)
über die modifizierte Knotenanalyse, ein Gleichungsblock liefert sie aus dem
Text seiner Gleichungen. Beides landet im selben System; deshalb lassen sich
Schaltung und Gleichungen beliebig koppeln.

Die Ableitungen für das Newton-Verfahren werden für alle Bauteile und alle
Ausdrücke automatisch und exakt gebildet. Sie müssen nirgends von Hand
angegeben werden.

\section{Die Eingabedatei in Kürze}

Eine Simulation ist vollständig durch eine Textdatei beschrieben. Sie ist
JSON mit Kommentaren (\code{//} und \code{/* */}):

\begin{json}
{
  "title": "Einweggleichrichter",

  "models": {
    "DGEN": { "type": "diode", "is": 2.5e-9, "n": 1.75, "cj0": "4p", "tt": "20n" }
  },

  "components": [
    { "type": "V", "name": "V1", "nodes": ["in", "0"],
      "wave": { "type": "sin", "amplitude": 5, "freq": "1k" } },
    { "type": "D", "name": "D1", "nodes": ["in", "out"], "model": "DGEN" },
    { "type": "R", "name": "R1", "nodes": ["out", "0"], "value": "10k" },
    { "type": "C", "name": "C1", "nodes": ["out", "0"], "value": "1u" }
  ],

  "analyses": [
    { "type": "op" },
    { "type": "tran", "tstop": "5m" }
  ],

  "outputs": ["v(in)", "v(out)", "i(V1)"]
}
\end{json}

Die Datei nennt Modelle, Bauteile mit ihren Knoten, die gewünschten
Analysen und die Größen, die ausgegeben werden sollen. Der Editor legt den
Schaltplan in derselben Datei unter dem Schlüssel \code{"schematic"} ab; der
Simulator überliest ihn. Eine im Editor gezeichnete Schaltung lässt sich
also jederzeit auch ohne Editor rechnen.

\section{Über dieses Handbuch}

\begin{description}[style=nextline,leftmargin=2em]
  \item[Teil I – Grundlagen] Installation und ein geführter erster Durchgang
    mit dem Editor.
  \item[Teil II – Der Editor] Alle Fenster, Werkzeuge und Tasten des
    Editors.
  \item[Teil III – Die Eingabedatei] Die vollständige Beschreibung des
    Formats: Zahlen, Ausdrücke, Bauteile, Modelle, Teilschaltungen,
    Analysen, Studien, Gleichungen, Ereignisse und Optionen. Dieser Teil
    gilt für Editor und Kommandozeile gleichermaßen.
  \item[Teil IV – Werkzeuge und Praxis] Kommandozeile, Python, SPICE-Import,
    Fehlersuche und ein Blick auf die Numerik.
  \item[Anhang] Eine vollständige Befehlsreferenz (Programme, alle
    Schlüssel der Eingabedatei mit Standardwerten, Ausdrücke, Befehle des
    Editors), Tastenkürzel, die mitgelieferten Beispiele, eine Kurzreferenz
    und Begriffe.
\end{description}

\paragraph{Schreibweisen.}
Text in \code{Schreibmaschinenschrift} steht so in Dateien oder auf der
Kommandozeile. Menüpunkte sind als \menue{Datei}\pfad\menue{Neu}
geschrieben, Tasten als \taste{R} oder \tasten{Strg}{R}.

\begin{hinweis}[Tasten auf dem Mac]
  Die Kürzel sind in diesem Handbuch in der Windows-Form angegeben. Auf dem
  Mac steht die Befehlstaste (\taste{cmd}) für \taste{Strg}, die Wahltaste
  (\taste{alt}) für \taste{Alt} und die Rückschritttaste für \taste{Entf}.
  Der Editor zeigt in Menüs, Tooltips und der Hinweiszeile jeweils die
  Namen des eigenen Systems.
\end{hinweis}
